\begin{table}[htbp]
\centering
\small
\setlength{\tabcolsep}{5pt}
\caption{Resumen de optimización intermedia y rol de cada variante.}
\label{tab:ch4_intermediate_optimization_summary}
\begin{tabular}{@{}>{\raggedright\arraybackslash}p{3.2cm}>{\raggedright\arraybackslash}p{8.1cm}@{}}
\toprule
Variante o etapa & Resultado experimental y uso posterior \\
\midrule
Búsqueda por familia & Identifica rangos plausibles de hiperparámetros para métodos candidatos sin fijar conclusiones definitivas. Se usa como evidencia intermedia para reducir el espacio de búsqueda. \\
Configuración global & Evalúa si una parametrización única mantiene desempeño razonable entre escenarios. Favorece reproducibilidad y evita reajuste por caso. \\
TRSS fijo & Variante estable con parámetros congelados antes de la comparación pareada. Es la lectura principal frente a MNE. \\
TRSS calibrado & Aproxima una selección más adaptativa, pero aún depende de reglas de calibración. Se usa como análisis auxiliar de sensibilidad. \\
Oráculo & Cota superior no desplegable al escoger la mejor variante observada. Sirve como contexto metodológico, no como método aplicable. \\
\bottomrule
\end{tabular}
\end{table}
